\subsection{PROJECTIVE SPACE STRUCTURE}

Given a set E and a field K, a (left) projective space structure on E with respect to the field K is defined by giving a non-empty set \(\Phi\) of bijections of subsets of the projective space \(\mathbf{P}(\mathbf{K}_{s}^{\mathrm{(E)}})\) onto E satisfying the following axioms:

\((\mathbf{EP}_{\mathrm{I}})\) The set of definition of every mapping \(f \in \Phi\) is a linear variety of \(\mathbf{P}(\mathbf{K}_s^{(\mathbf{E})})\) .

\((\mathrm{EP}_{\Pi})\) For every ordered pair of elements \(f, g\) of \(\Phi\) defined respectively on the linear varieties \(\mathbf{P}(\mathbf{V})\) and \(\mathbf{P}(\mathbf{W})\) , the bijection \(h = \overline{g}^{-1} \circ f\) of \(\mathbf{P}(\mathbf{V})\) onto \(\mathbf{P}(\mathbf{W})\) is a projective mapping.

\((\mathbf{EP}_{\mathrm{III}})\) Conversely, if \(f \in \Phi\) is defined on the linear variety \(\mathbf{P}(\mathbf{V})\) and \(h\) is a bijective projective mapping of \(\mathbf{P}(\mathbf{V})\) onto a linear variety \(\mathbf{P}(\mathbf{W}) \subset \mathbf{P}(\mathbf{K}_s^{\mathbf{(E)}})\) , then \(f \circ h^{-1} \in \Phi\) .

Let E be a set, \((\mathbf{V}_{\lambda})_{\lambda \in L}\) a family of vector spaces over K and suppose given for each \(\lambda \in L\) a bijection \(f_{\lambda}\) of \(\mathbf{P}(\mathbf{V}_{\lambda})\) onto E such that, for every ordered pair of indices \(\lambda, \mu, f_{\lambda} \circ f_{\mu}\) is a projective mapping of \(\mathbf{P}(\mathbf{V}_{\mu})\) onto \(\mathbf{P}(\mathbf{V}_{\lambda})\) . Then we can define on E a projective space structure with respect to K as follows: let \((e_{t})_{t \in I}\) be a basis of a space \(V_{\lambda}\) and write \(a_{t} = f_{\lambda}(\pi(e_{t}))\) ; let \(b_{t}\) be the element of index \(a_{t}\) in the canonical basis of \(\mathbf{K}_{s}^{(\mathbf{E})}\) (§1, no. 11). The relation \(t \neq k\) implies \(b_{t} \neq b_{k}\) because of the hypothesis that \(f_{\lambda}\) is bijective; hence the \(b_{t}\) form a basis of a vector subspace \(W_{0}\) of \(K_{s}^{(\mathbf{E})}\) and there therefore exists a bijective projective mapping h of \(\mathbf{P}(\mathbf{W}_{0})\) onto \(\mathbf{P}(\mathbf{V}_{\lambda})\) such that \(h(\pi(b_{t})) = \pi(e_{t})\) for all \(t \in I\) . If \(\Phi\) is taken to be the set of all bijective projective mappings \(f_{\lambda} \circ h \circ g^{-1}\) , where g runs through the set of all bijective projective mappings \(\mathbf{P}(\mathbf{W}) \subset \mathbf{P}(\mathbf{K}_{s}^{(\mathbf{E})})\) , it is immediately verified that \(\Phi\) satisfies axioms \((\mathrm{EP}_{\mathrm{I}}), (\mathrm{EP}_{\mathrm{II}})\) and \((\mathrm{EP}_{\mathrm{III}})\) . It is moreover immediate that \(\Phi\) depends neither on the choice of index \(\lambda \in L\) , nor on the choice of basis \((e_{t})\) in \(V_{\lambda}\) , nor on the choice of h.

In particular (taking L to consist of a single element), every projective space \(\mathbf{P}(\mathbf{V})\) derived from a vector space V (no. 5, Definition 4) thus has a well determined ``projective space structure'' in the sense of the definition given in this no. Hence any set with a projective space structure can be called a projective space.

\[
f \in \Phi
\]

\[
\mathbf {P} (\mathbf {V}) \subset \mathbf {P} \left(\mathrm{K} _ {\mathrm{s}} ^ {(\mathrm{E})}\right)
\]

\(f(\mathbf{M})\) is a linear variety in \(\mathbf{P}(\mathbf{V})\) in the sense of no. 7 (this property then holds for all \(f \in \Phi\) ). It follows from the above that every linear variety in a projective space has canonically a projective space structure.

A projective space \(\mathbf{E}\) is said to be of dimension \(n\) if, for all \(f\in\Phi,f(\mathbf{E})\) is a linear variety of dimension \(n\) (it suffices that this hold for one mapping \(f\in\Phi\) ).

